

\section{Electroweak Corrections in the Standard Model}

\subsection{Contribution (Results of top-Yukawa induced corrections):} 

\textit{Authors:} Margarete Mühlleitner, Johannes Schlenk, Michael Spira (2207.02524) \\[0.3cm]
The first calculation of partial electroweak corrections to Higgs-pair production via the dominant gluon-fusion process has been performed by investigating the corrections induced by the top-Yukawa coupling \cite{Muhlleitner:2022ijf}. In order to set-up a complete and consistent framework, these corrections have been obtained in the gaugeless limit, where all electroweak gauge couplings are set to zero, but the (massless) Goldstone contributions are taken into account. This approach ensures the preservation of the $SU(2)$ symmetry of the scalar Higgs doublet. While the top-induced corrections to the trilinear Higgs vertex are treated with the full top-mass dependence, the triangle and box diagrams involving single- and double-Higgs couplings to gluons are treated in the heavy-top limit (HTL). The corresponding diagrams are displayed in Fig.~\ref{fg:dia_htl}, where the first two represent those that are treated in the HTL. For these two diagrams, we have used the effective Lagrangian \cite{Djouadi:1994ge,Muhlleitner:2022ijf}
\begin{figure}[hbt]
%\vspace*{-1cm}
\begin{center}
\includegraphics[width=0.6\textwidth]{figures/dia_htl}
%\vspace*{-7cm}
\caption{\label{fg:dia_htl} \it Generic diagrams contributing to the top-Yukawa induced electroweak corrections to the $gg\to HH$ process. The first two diagrams have been treated in the HTL.}
\end{center}
\end{figure}
\begin{equation}
{\cal L}_{eff} = \frac{\alpha_s}{12\pi} G^{a\mu\nu} G^a_{\mu\nu} \left\{
(1+\delta_1) \frac{H}{v} + ( 1 + \eta_1) \frac{H^2}{2v^2} + {\cal
O}(H^3) \right\}
\label{eq:leff}
\end{equation}
where
\begin{eqnarray}
\delta_1 & = & \frac{x_t}{2} + {\cal O}(x_t^2) \,, \qquad\qquad
\eta_1   = 4 x_t + {\cal O}(x_t^2) \,.
\label{eq:leff_coeff}
\end{eqnarray}
with $x_t = G_F m_t^2/(8\sqrt{2}\pi^2)$, $G_F$ denotes the Fermi constant, $v$ the vacuum expectation value (vev) of the Higgs field and $m_t$ the top mass. This effective Lagrangian describes the electroweak corrections induced by $x_t$ to the $Hgg$ and $HHgg$ vertices in the HTL. We would like to point out explicitly that the square root of the wave-function counterterm of the external Higgs boson(s) is already taken into account in this effective Lagrangian.

In Fig.~\ref{fg:dia_htl} the tadpole diagrams are displayed explicitly. We have used both, the Fleischer--Jegerlehner and the conventional scheme of omitting the tadpole diagrams, but keep them in the corresponding counterterm of the trilinear Higgs coupling. We have obtained identical results in the final sum of the radiative corrections, once the vev $v$ is expressed in terms of the Fermi constant $G_F$.

Because the full top-mass dependence of the one-loop times one-loop diagrams describing the radiatively corrected triple Higgs vertex (all diagrams of Fig.~\ref{fg:dia_htl} except the first two) have been kept, the use of the effective trilinear Higgs coupling (derived from the effective Coleman--Weinberg potential \cite{Coleman:1973jx,Weinberg:1973ua,Jackiw:1974cv})
\begin{equation}
\lambda_{HHH} = 3 \frac{M_H^2}{v} \, , \qquad\qquad \lambda_{HHH}^{eff} = \lambda_{HHH} - \frac{3 m_t^4}{\pi^2 v^3} \approx 0.91 \times \lambda_{HHH}
\end{equation}
to accommodate the dominant part of the radiative corrections can be tested. The result is that process-dependent, finite-momentum-dependent corrections are of the same size as the corrections above, growing with $m_t^4$,
\begin{eqnarray}
\sigma & = & K_{elw} \times \sigma_{LO} \nonumber \\
K_{elw} & \approx & 1.002 \hspace*{1.0cm} \mbox{($\lambda_{HHH}$)} \nonumber \\
K^{eff}_{elw} & \approx & 0.938 \hspace*{1cm} \mbox{($\lambda^{eff}_{HHH}$)}
\end{eqnarray}
Since the effective trilinear coupling generates about 6\% electroweak corrections to the total cross section artificially and larger corrections to the distributions, using the LO-like trilinear coupling $\lambda_{HHH}$ is favored \cite{Muhlleitner:2022ijf}. The top-Yukawa induced electroweak corrections are small except in the region close to the production threshold due to the complete cancellation of the leading top-mass terms in the LO matrix element; which the corrections are normalized to.

\subsection{Contribution (Top-Yukawa and light-quark induced corrections with full mass dependence):}

\textit{Authors:} Arunima Bhattacharya, Francisco Campanario, Sauro Carlotti, Jamie Chang, Javier Mazzitelli, Margarete Mühlleitner, Jonathan Ronca, Michael Spira (in preparation)
\subsubsection{Top-Yukawa induced electroweak corrections}
The analysis of the top-Yukawa induced electroweak corrections discussed above has been extended to include the full top-quark mass dependence. Specifically, the first two diagrams of Fig.~\ref{fg:dia_htl}, previously evaluated in the heavy-top limit, have been resolved into the corresponding two-loop diagrams involving top, bottom, Higgs, and Goldstone propagators. To preserve the SU(2) symmetry of the Higgs sector, this work has been performed in the gaugeless limit, enabling a gauge-invariant definition of the contributions associated with the top-Yukawa coupling. In this limit, the Goldstone bosons are massless but do not generate infrared singularities. The method used for this calculation involves applying the projectors to the two form factors (in $D=4-2\epsilon$ dimensions), followed by a diagram-by-diagram Feynman parametrization. Since no tensor reduction has been used, the resulting Feynman integrands remain quite compact. The singularities were extracted by appropriate end-point subtractions. Moreover, this allowed us to keep, alongside the top mass, the bottom and the Higgs masses as free symbolic parameters. The virtual thresholds have been treated by introducing complex masses of the virtual propagators,
\begin{equation}
m_{t/b}^2 \to m_{t/b}^2 (1-i\bar\epsilon) \, , \qquad M_H^2 \to M_H^2 (1-i\bar\epsilon)
\end{equation}
with the regulator $\bar\epsilon$ to be sent to 0. The latter was achieved by applying a Richardson extrapolation \cite{Richardson} to a series of different values of $\bar\epsilon$. The minimal value of $\bar\epsilon$ used for numerical integrations varied between $10^{-8}$ and 0.01 depending on the diagram and the value of the invariant Higgs pair mass $Q=M_{HH}$. The top mass has been renormalized on-shell, and the vev according to the proper definition of the Fermi constant in muon decay.

\begin{figure}[hbt]
\begin{center}
\includegraphics[width=0.6\textwidth]{figures/dia_top}
\caption{\label{fg:dia_top} \it Sample two-loop triangle and box diagrams of the top-Yukawa-induced electroweak corrections to Higgs-boson pair production involving Higgs $h$ and Goldstone $G^0,G^\pm$ exchanges. The bottom propagators only contribute to the diagrams with charged Goldstone exchange.}
\end{center}
\end{figure}
\begin{figure}[hbt]
\begin{center}
\includegraphics[width=0.5\textwidth]{figures/d1tria}
\caption{\label{fg:delta_1} \it The full top-mass dependence of the coefficient $\delta_1$ of Eq.~(\ref{eq:leff_coeff}) as a function of the invariant Higgs pair mass. The dotted line shows the purely real HTL result. The sizable imaginary part arises from the diagrams with charged Goldstone exchange that include diagrams with $b\bar b$ thresholds. The points on the curve indicate the $M_{HH}$ values for which $\delta_1$ has been computed. The numerical errors are negligible and thus not shown.}
\end{center}
\end{figure}
The result of the two-loop triangle diagrams, see Fig.~\ref{fg:dia_top}, which can be directly compared to the coefficient $\delta_1$ in Eq.~(\ref{eq:leff_coeff}) is shown in Fig.~\ref{fg:delta_1}. As a cross-check, the HTL result of Eq.~(\ref{eq:leff_coeff}) has been reproduced within uncertainties by artificially increasing the value of the top mass. The results indicate differences between the HTL and the full result at the (sub)percent level for the distribution in the Higgs pair mass.

The calculation of the two-loop box diagrams, see Fig.~\ref{fg:dia_top}, is more involved, since the additional phase space dependence does not allow for an immediate comparison with the HTL. Applying the same method as for the two-loop triangle diagrams and including the phase space integration in the numerical analysis, the result of the two-loop renormalized box diagrams plus the two-loop triangle diagrams and the previous one-loop times one-loop diagrams related to the radiative corrections to the trilinear Higgs vertex is shown in Fig.~\ref{fg:del_top}, where the relative corrections are defined as
\begin{equation}
\sigma (gg\to HH) = \sigma_{LO} (gg\to HH)~( 1 + \delta_{top-Yukawa} + \delta_{light-quarks} )
\end{equation}
with the light-quark contributions $\delta_{light-quarks}$ discussed in Section \ref{sc:light-quarks}. The total corrections are about --5\% for intermediate values of $M_{HH}$ and are smaller for large values, while they are large close to the production threshold due to the large cancellation in the LO matrix element, to which the corrections are normalized. At large values of $M_{HH}$, the corrections develop a visible slope. The right plot of Fig.~\ref{fg:del_top} shows details about the $t\bar t$ threshold that develops a sign-changing interference behavior with the LO matrix element, with a maximum below the $t\bar t$ threshold.
\begin{figure}[hbt]
\vspace*{-3cm}
\includegraphics[width=0.5\textwidth]{figures/del_top}
\vspace*{-10cm}
\hspace*{-0.5cm}
\includegraphics[width=0.5\textwidth]{figures/del_top.detail}
\vspace*{6.7cm}
\caption{\label{fg:del_top} \it The full relative top-Yukawa induced electroweak corrections to the $gg \to HH$ process. The points on the curve indicate the $M_{HH}$ values for which the corrections have been computed. The right plot shows a magnified view on the structure of the $t\bar t$ threshold. The numerical error bars have been added to each of the points. The vertical line indicates the position of the $t\bar t$ threshold.}
\end{figure}

\subsubsection{Light-quark loops} \label{sc:light-quarks}
The contributions of closed light-quark loops to the $gg\to HH$ process have been calculated as well in the full electroweak gauge theory. The two-loop triangle diagrams (see the first diagram of Fig.~\ref{fg:dia_lq}) can be adopted from the single-Higgs calculation \cite{Aglietti:2004nj,Aglietti:2006yd} and translated to the two-loop box diagrams with a 4-point vertex (see the second diagram of Fig.~\ref{fg:dia_lq}). We have included four light flavors in the diagrams involving charged $W$ and Goldstone boson exchange, while in diagrams with neutral $Z$ and Goldstone boson exchange five flavors have been taken into account. The bottle-neck of this calculation is the evaluation of the genuine two-loop box diagrams (see the last two diagrams of Fig.~\ref{fg:dia_lq}). The numerical stability of these diagrams is quite demanding, since strong cancellations emerge between all diagrams, which requires the form factors of the individual diagrams to be numerically determined with very significant precision. The numerical method applied to these diagrams is the same as before, i.e.~projection on the two form factors in $D=4-2\epsilon$ dimensions, no tensor reduction, Feynman parametrization and analytical continuation of all propagator masses including the $W$ and $Z$ boson masses, $M_{W/Z}^2 \to M_{W/Z}^2 (1-i\bar\epsilon)$. To isolate the singularities, multiple end-point subtractions need to be performed. In order to reach the narrow-width limit $\bar\epsilon\to 0$, a Richardson extrapolation was performed with a minimal regulator down to $\bar\epsilon = 0.025$ depending on $M_{HH}$.
\begin{figure}[hbt]
\vspace*{-1.5cm}
\begin{center}
\includegraphics[width=0.7\textwidth]{figures/dia_lq}
\vspace*{-12cm}
\caption{\label{fg:dia_lq} \it Sample two-loop box diagrams of the light-quark-induced electroweak corrections to Higgs-boson pair production. Shown are diagrams of the three classes -- triangle diagrams, four-point diagrams and genuine planar/non-planar double boxes.}
\end{center}
\end{figure}
\begin{figure}[hbt]
\vspace*{-3.5cm}
\begin{center}
\includegraphics[width=0.5\textwidth]{figures/delta_lq}
\vspace*{-3.5cm}
\caption{\label{fg:delta_lq} \it The full relative light-quark loop induced electroweak corrections to the $gg \to HH$ process. The points on the curve together with their numerical error bars indicate the $M_{HH}$ values for which the corrections have been computed.}
\end{center}
\end{figure}

The total sum of all light-quark loop diagrams is finite, and no renormalization is required. The result is a small correction at the sub-percent level for large values of $M_{HH}$, see Fig.~\ref{fg:delta_lq}, while the corrections are larger close to the production threshold due to the cancellation of the leading top-mass terms of the LO matrix element, to which the results are normalized. In total, the light-quark loops do not play a dominant role in Higgs pair production via gluon fusion in contrast to single-Higgs production \cite{Actis:2008ts,Actis:2008ug}. The results agree with the results of Ref.~\cite{Bonetti:2025vfd} within the numerical uncertainties; see Section \ref{sc:bonetti}.

\subsection{Contributions (Top-Yukawa and Higgs self-coupling corrections with full top quark mass dependence)}

\textit{Authors:} Gudrun Heinrich, Stephen Jones, Matthias Kerner, Thomas Stone, Augustin Vestner (2407.04653) 

\begin{figure}[ht]
    \centering
    \begin{subfigure}{0.25\textwidth}
    \centering
    \fig{figures/figs_240704653/NLOtritpropHbubH}
    \end{subfigure}
    \hspace{5em}  
    %
    \begin{subfigure}{0.25\textwidth}
    \centering
    \fig{figures/figs_240704653/NLOboxHvertcorr}
    \end{subfigure}
    \hspace{5em}  
    %
    \begin{subfigure}{0.25\textwidth}
    \centering
    \fig{figures/figs_240704653/NLOboxtboxt}
    \end{subfigure}
    %
    \caption{Example diagrams contributing to the different coupling structures on which the bare two-loop amplitude is separated.}
    \label{fig:240704653_examplediags}
\end{figure}

In Ref.~\cite{Heinrich:2024dnz}, fully differential results for the Yukawa- and Higgs self-coupling-type electroweak corrections to di-Higgs production in gluon fusion are presented.
The results retain the full dependence on the mass of the top quark and the Higgs boson.
At the amplitude level, the results are further separated on the individual coupling structures, allowing the bare amplitude to be computed with arbitrary top-Yukawa and Higgs trilinear and quartic couplings.

The precise class of corrections which are computed is defined via a Yukawa model with only an up-type quark (the top quark) and a scalar field (the Higgs boson), example Feynman diagrams are shown in Fig.~\ref{fig:240704653_examplediags}. 
This model selects a well-defined subset of diagrams which can be obtained from considering the SM in the unitary gauge and then setting $(g, g^\prime) \rightarrow (0,0)$, which removes the electroweak gauge bosons (and their associated ghost fields).
The electroweak input-parameter scheme $\left\{ M_Z=0, M_W=0, G_F\right\}$ + $\left\{m_t,m_H\right\}$ is used, where all masses are specified in the on-shell scheme.
Results are presented with the Higgs vev fixed in the $G_\mu$ scheme (with $M_Z=0$, $M_W=0$), however, the intermediate results are sufficiently general to allow another renormalisation scheme to be adopted.
Tadpole contributions are treated within the Fleischer-Jegerlehner tadpole scheme (FJTS)~\cite{Fleischer:1980ub}.

The calculation is performed by projecting the bare amplitude onto the two form factors described above in $D=4-2\epsilon$ dimensions.
The unreduced amplitude is obtained using two separate calculations based on either \texttt{alibrary}~\cite{alibrary}, a \texttt{Mathematica} and \texttt{Form}~\cite{Kuipers:2013pba} package for computing multi-loop amplitudes, or \texttt{Reduze 2}~\cite{vonManteuffel:2012np}, allowing for a detailed cross-check.
The form factors are then reduced fully symbolically, retaining the dependence on the Mandelstam invariants ($s, t$), the masses ($m_H, m_t$) and the space-time dimension $D$, to a basis of 494 finite master integrals using the public programs \texttt{Kira}~\cite{Klappert:2020nbg}, \texttt{Ratracer}~\cite{Magerya:2022hvj}, and \texttt{Firefly}~\cite{Klappert:2019emp,Klappert:2020aqs}.
The master integrals appearing in the reduced form factors are then evaluated numerically for each phase-space point using \texttt{pySecDec}~\cite{Borowka:2017idc,Borowka:2018goh,Heinrich:2021dbf,Heinrich:2023til}, the evaluation of the master integrals is additionally cross-checked at several physical points using \texttt{DiffExp}~\cite{Hidding:2020ytt,Moriello:2019yhu}.

\begin{figure}[h!]
\centering
\begin{subfigure}{0.4\textwidth}
\centering\includegraphics[width=1.0\textwidth]{figures/figs_240704653/mhh_full}
\end{subfigure}
%
\qquad
%
\begin{subfigure}{0.4\textwidth}
\centering\includegraphics[width=1.0\textwidth]{figures/figs_240704653/pth_full}
\end{subfigure}
\caption{Invariant mass and transverse momentum distributions for di-Higgs boson production at LO and $\mathrm{NLO}^\mathrm{EW}$ including only the Yukawa and self-coupling type corrections.}
\label{fig:240704653_dists}
\end{figure}

In Fig.~\ref{fig:240704653_dists}, results for the di-Higgs invariant mass and $p_{T,H}$~distribution are shown.
The results presented here are obtained using the \texttt{PDF4LHC21\_40}~\cite{PDF4LHCWorkingGroup:2022cjn} distribution functions interfaced via \texttt{LHAPDF}~\cite{Buckley:2014ana} and setting the factorisation and renormalisation scale to $\mu_r = \mu_f = m_{HH}/2$. The masses of the Higgs boson and top quark are set to $m_H = 125~\mathrm{GeV}$, $m_t = \sqrt{23/12}\,m_H =173.055~\mathrm{GeV}$, respectively, and $G_F = 1.1663787 \cdot 10^{-5} \ \mathrm{GeV}^{-2}$, corresponding to $v=246.22~\mathrm{GeV}$.
For the subset of the electroweak correction considered, very large shape distortions (around $30\%$ for the binning selected) are observed close to threshold for the invariant mass distribution, with much smaller corrections above 400 GeV (up to 1 TeV).
The corrections in the $p_{T,H}$ distribution are less localised, with distortions of up to $5\%$ present across the spectrum.
Overall, the top-Yukawa and Higgs self-coupling corrections amount to a $1\%$ enhancement of the total cross-section.

An important observation concerns the renormalisability of models in which the top-Yukawa, Higgs boson trilinear and quartic couplings are varied from their SM values.
In particular, it is observed that the $1/\epsilon$ pole structure of the vev counter-term obtained by demanding the finiteness of the Yukawa or self-coupling vertices differs unless they are each set to their SM value.
This implies that the $\kappa$-framework used to extract the modifier of the Higgs trilinear coupling, $\kappa_3$, can not directly be used when including the electroweak corrections to this process.
Furthermore, even though the quartic Higgs boson self-coupling first enters at $\mathrm{NLO}_\mathrm{EW}$, the subset of diagrams in which it appears are divergent prior to UV renormalisation, implying that also the $\kappa_4$ coupling can not naively be varied.

\subsection{Contribution (Analytic results through expansions):}

\textit{Authors:} Joshua Davies, Kay Schönwald, Matthias Steinhauser, Hantian Zhang (2501.17920)

There are a number of calculations where analytic results
have been obtained for the NLO electroweak corrections.

The full contribution including all sectors of the SM has been considered in
Ref.~\cite{Davies:2023npk} in an expansion in the inverse top quark mass. Five
expansion terms in $1/m_t^2$ have been computed. Although the radius of convergence is
limited, this result serves as an important benchmark for other calculations,
see, e.g., Ref.~\cite{Bi:2023bnq} and the discussion below.

In Ref.~\cite{Zhang:2024rix} exact analytic results for the factorizable
contributions have been obtained. No approximation has been
applied. Computer-readable results for the form factors are provided which are
expressed in terms of standard $B_0$ and $C_0$ functions.

Analytic high-energy results for the subset of Feynman diagrams involving top quark
Yukawa and Higgs boson self-energy couplings have been computed in
Refs.~\cite{Davies:2022ram,Davies:2025wke}.  Such diagrams involve as
dimensionful quantities the Mandelstam variables $s$ and $t$, the top quark
mass and the Higgs boson mass.  In the calculation the final-state Higgs mass
and the one in the propagators inside the loop diagrams are treated separately. This
allows us to perform as a first step an expansion in the final-state mass and
subsequently an expansion in $\delta = 1-m_H^{\rm int}/m_t$ with $m_H^{\rm int}$ being the
internal propagator mass. Both expansions are simple Taylor expansions.
Afterwards, we realize the high-energy limit through an expansion for small
$m_t$ which requires the application of a nontrivial asymptotic expansion.
%
Technical details of the master integral calculations in the high-energy limit are given in Refs.~\cite{Zhang:2024fcu,Davies:2022ram,Davies:2025wke},
together with a publicly available tool \texttt{AsyInt}~\cite{Zhang:2024fcu}.
%
We typically compute about 100 expansion terms in $m_t$ and subsequently apply
a Pad\'e procedure that provides, for each phase-space point, a central value
and an uncertainty. It has been shown in Refs.~\cite{Davies:2022ram, Davies:2023vmj} that this
approach leads to precise results down to
fairly low values of the Higgs transverse momentum $p_T$ in the case of QCD and leading Yukawa corrections.

As an illustration we show in Fig.~\ref{fig::diff_to_SD_yt3lam1} the
comparison of the real part of the box contribution to form factor $F_1$
originating from diagrams with three top quark Yukawa couplings and one Higgs
boson self-coupling (see Fig.~\ref{fig::diags} for sample diagrams) to the
numerical results from Ref.~\cite{Heinrich:2024dnz}. After including expansion
terms up to order $m_H^4\delta^3$ one observes agreement at the percent level.

Note that the analytic results provide full flexibility concerning the
change of renormalization schemes and the variation of input parameters.

\begin{figure}[t]
  \centering
  \begin{tabular}{ccc}
  \includegraphics[width=0.32\textwidth]{figures/figs_250117920/diff_to_SD_F1bare_ep0_t-plus_mhs2del3_yt3lam1_pTcut300_re_sqrts_FF.pdf}
  \includegraphics[width=0.32\textwidth]{figures/figs_250117920/diff_to_SD_F1bare_ep0_t-plus_mhs2del3_yt3lam1_pTcut300_re_sqrts.pdf}
  \includegraphics[width=0.32\textwidth]{figures/figs_250117920/diff_to_SD_F1bare_ep0_t-plus_mhs2del3_yt3lam1_pTcut300_re_sqrts_rel.pdf}
   \end{tabular}
   \caption{\label{fig::diff_to_SD_yt3lam1} Real part of the form factor (left
     column) from Feynman diagrams like those shown in Fig.~\ref{fig::diags}.
     The middle and right columns show the absolute and relative difference compared
     to Ref.~\cite{Heinrich:2024dnz}.  }
\end{figure}


\begin{figure}[t]
  \centering
  \begin{tabular}{c}
    \includegraphics[width=0.3\textwidth]{figures/figs_250117920/gghh_c.pdf} 
    \includegraphics[width=0.28\textwidth]{figures/figs_250117920/gghh_e.pdf} 
  \end{tabular}
  \caption{\label{fig::diags}Sample Feynman diagrams. Solid, dashed and curly
    lines represent top quarks, Higgs bosons and gluons, respectively.}
\end{figure}

\subsection{Contribution:} \label{sc:bonetti}

\textit{Authors:} Marco Bonetti, Philipp Rendler, William J. Torres Bobadilla (2503.16620)

~



We provide fully symbolical expressions of the amplitude for double Higgs boson production in gluon fusion mediated by light quarks~\cite{Bonetti:2025vfd}.


\begin{figure}
\centering
    \subfloat[]{{\includegraphics[height=0.11\textheight]{Bonetti_Rendler_Torres_Bobadilla/Images/gg-HHH}}\label{fig:VVHV}}
	\qquad\qquad
    \subfloat[]{{\includegraphics[height=0.11\textheight]{Bonetti_Rendler_Torres_Bobadilla/Images/gg-HHVV.pdf}}\label{fig:HHVV}}
	\qquad\qquad
    \subfloat[]{{\includegraphics[height=0.11\textheight]{Bonetti_Rendler_Torres_Bobadilla/Images/gg-VVV}}\label{fig:VVV}}
	\caption{Representatives for the light-quark contributions.}
    \label{fig:amplgg}
\end{figure}


This mode of production appears for the first time at two loops: the gluons couple to a light-quark box, to which either two $W$ or two $Z$ bosons are attached, subsequently generating the Higgs boson(s). We divide the relevant diagrams into three blocks: one-particle reducible diagrams containing one vector-vector-Higgs and one triple-Higgs vertex (dubbed $VVh$, see Fig.~\ref{fig:VVHV}); diagrams containing a vector-vector-Higgs-Higgs vertex ($VVhh$, Fig.~\ref{fig:HHVV}); and diagrams containing two vector-vector-Higgs vertices ($VVV$, Fig.~\ref{fig:VVV}). Each one of these blocks is further divided into diagrams containing either $W$ or $Z$ bosons only. Each one of the aforementioned terms is explicitly gauge-invariant and UV- and IR-finite.

No axial terms are generated by the fermion loop thanks to charge-parity conservation. The amplitude can then be decomposed as a linear combination of the standard tensors $T_1^{\mu\nu}$ and $T_2^{\mu\nu}$ already used for double Higgs production when considering QCD corrections~\cite{Borowka:2016ehy}. We consider the first two generations of light quarks coupling to the $W$ boson and five light quarks coupling to the $Z$ boson (all but the top quark).

We extract the two form factors and reduce the set of scalar two-loop Feynman integrals appearing in the amplitude to master integrals. We derive differential equations for specific linear combinations of the master integrals w.r.t.\ the Mandelstam variables and the masses, such that their solution can be written in terms of iterated integrals over logarithmic kernels. We improve on this choice of basis by applying a series of rotations to eliminate redundant integration kernels. We fix the integration constants by matching the expressions of the master integrals to their large-$m_{W,Z}$ expansion.

We employ the results to obtain analytic expressions for the form factors, and we further identify independent functions therein for which we build dedicated differential equations for optimized numerical evaluation, which we solve in terms of generalized series expansions, evolving from the large-mass limit to arbitrary points in the physical region of the phase space. We provide a code implementation for the numerical evaluation of our functions based on the \texttt{Mathematica} package \texttt{DiffExp}~\cite{Hidding:2020ytt}. Grids for the light-quark amplitude, as well as for their interference with the LO amplitude have been implemented in the POWHEG-BOX framework or are available upon request.

We observe that the interference between the LO and the light-quarks first form factor represents the main source of contributions except for the high-energy tail, where the second form factor becomes relevant, although very small overall. We notice large cancellations when considering the interference between LO triangle and box diagrams with each of the $VVH$, $VVHH$, and $VVV$ blocks of diagrams, suggesting that the light-quark terms are highly sensitive to variations of the tri-linear Higgs coupling.

The results have been checked against the independent evaluation~\cite{Bhattacharya_TBA}.

~

We present the impact of Yukawa~\cite{Heinrich:2024dnz} and light-quark corrections~\cite{BoHeReToBa_TBA} to the Higgs-pair invariant mass distribution and the transverse momentum distribution in Fig.~\ref{fig:mHH} and \ref{fig:pTH}, respectively.


\begin{figure}
\centering
    \subfloat[]{{\includegraphics[width=0.48\textwidth]{Bonetti_Rendler_Torres_Bobadilla/Images/16_mhh2_Yukawa_s136}}\label{fig:mHH}}
	\quad
    \subfloat[]{{\includegraphics[width=0.48\textwidth]{Bonetti_Rendler_Torres_Bobadilla/Images/17_pt_Yukawa_s136}}\label{fig:pTH}}
	\caption{Light-quark and Yukawa contributions to the Higgs pair invariant mass and Higgs transverse momentum distributions.}
    \label{fig:obs}
\end{figure}

In the Higgs-pair invariant mass distribution both classes of contributions affect mainly the low-energy region, where the Yukawa terms produce an enhancement of almost $+40\%$, while the light-quarks have a negative effect of more than $-15\%$. Both contributions vanish at high energies. While the light quark have a constant negative effect on the distribution, the Yukawa contributions become slightly negative only between $400~\textup{GeV}$ and $800~\textup{GeV}$. Their combined effect produces an increase of more than $+20\%$ at the production threshold, rapidly approaching zero as the energy grows.

A similar behaviour for light quarks is observed in the transverse momentum distribution, where they have a negative effect at low energies of around $-4\%$, slowly approaching zero for increasing energies. The Yukawa contributions have an oscillatory pattern: they raise to $+5\%$ between $50\,\textup{GeV}$ and $150\,\textup{GeV}$, become negative between $200\,\textup{GeV}$ and $400\,\textup{GeV}$, and slowly rise again at higher energies.






\subsection{Contribution (Full EW in SM):}
\label{sec:hhew}

\textit{Authors:} Huan-Yu Bi, Li-Hong Huang, Rui-Jun Huang, Yan-Qing Ma, Huai-Min Yu (2311.16963)

The ampltiude for $g(p_1)g(p_2) \to H(p_3)H(p_4)$ can be expressed as ${\cal M}=\varepsilon_{1\mu}\varepsilon_{2\nu}{\cal M}^{\mu\nu}$, where $\varepsilon_{i\mu}$ denotes gluon polarizaiton vector. Due to the current conservation relations $p_{1\mu}{\cal M}^{\mu\nu}=p_{2\nu}{\cal M}^{\mu\nu}=0$, we have the decomposition
\begin{align}
{\cal M}^{\mu\nu}=F_1T_1 ^{\mu\nu}+F_2T_2 ^{\mu\nu}+\Delta_0^{\mu\nu}+\Delta_5^{\mu\nu},
\end{align}
where $\Delta_5^{\mu\nu}$ contains only one Levi-Civita tensor and appears from two-loop order onward, and $\Delta_0^{\mu\nu}$ depends on either $p_1^{\mu}$ or $p_2^{\nu}$. The later two tensors will vanish for NLO EW calculations when multiplied by the  LO amplitude. The physical tensors can be expressed as \cite{Plehn:1996wb}
\begin{align}
&T_1 ^{\mu\nu}= g^{\mu\nu}-\frac{p_1^{\nu}\,p_2^{\mu}}{p_1\cdot  p_2} \;,\label{eq:Ttensors} \\
&T_2 ^{\mu\nu}= g^{\mu\nu}+\frac{1}{p_{T}^{2}\,(p_1\cdot p_2)}\big[ 2\,(p_1\cdot p_2) \,p_3^{\nu}\,p_3^{\mu}\nonumber  \\
&~~ - 2\,(p_1\cdot p_3) \,p_3^{\nu}\,p_2^{\mu} -2\,(p_2\cdot p_3) p_3^{\mu}\,p_1^{\nu}+ m_H^2 \,p_1^{\nu}\,p_2^{\mu} \big] \;, 
\end{align}  
where $p_{T}=\sqrt{(\hat{t}\hat{u} - m_H^4)/\hat{s}}$ represents the transverse momentum of one of the Higgs with 
\begin{eqnarray}
&&\hat{s}=(p_1+p_2)^2,~\hat{t}=(p_1-p_3)^2,~\hat{u}=(p_2-p_3)^2.\label{st}
\end{eqnarray} 
The gauge-invariant form factor $F_{1,2}$ can be extracted by defining the projectors $P_{1,2}^{\mu\nu}$ in terms of $T_{1,2}^{\mu\nu}$
\begin{eqnarray}
P_1^{\mu\nu} &=&+\frac{1}{4}\,\frac{D-2}{D-3} \,T_1^{\mu\nu}
-\frac{1}{4}\,\frac{D-4}{D-3} \,T_2^{\mu\nu}\;,\label{eq:proj1}\\
P_2^{\mu\nu}&=& -\frac{1}{4}\,\frac{D-4}{D-3} \,T_1^{\mu\nu}
+\frac{1}{4}\,\frac{D-2}{D-3} \,T_2^{\mu\nu}\;,\label{eq:proj2}
\end{eqnarray}   
such that 
\begin{eqnarray}
F_1=P_1^{\mu\nu} {\cal M}_{\mu\nu}\;,~ F_2=P_2^{\mu\nu} {\cal M}_{\mu\nu}\;.
\end{eqnarray}
Crucially, real-emission processes $gg \to HH\gamma$ are forbidden by the Furry theorem, restricting NLO EW corrections to purely virtual contributions at the two-loop level.

In Ref.~\cite{Bi:2023bnq}, it was found that there are 116 integral families for the NLO EW corrections. The loop integrals in each family are reduced to master integrals using the {\tt Blade} package \cite{Blade}. It utilizes {\tt FiniteFlow} \cite{Peraro:2019svx} to solve integration-by-parts (IBP) relations ~\cite{Chetyrkin:1981qh}, incorporating a block-triangular form \cite{Liu:2018dmc,Guan:2019bcx} to enhance computational efficiency. Dimensional regularization ($D=4-2\epsilon$) with  $\epsilon=\pm1/1000$ is chosen, which eliminates the need for Laurent expansions in $\epsilon$ for cross sections and thus reduces resource demands \cite{Liu:2022chg,Liu:2022mfb}. Master integrals are evaluated by numerically solving systems of differential equations \cite{Kotikov:1990kg,Remiddi:1997ny,Caffo:2008aw,Czakon:2008zk,Lee:2014ioa,Moriello:2019yhu,Hidding:2020ytt,Armadillo:2022ugh} in $\hat{s}$ and $\hat{t}$,  with boundary conditions computed by {\tt AMFlow} package \cite{Liu:2022chg} implementing the auxiliary mass flow method \cite{Liu:2022mfb,Liu:2021wks,Liu:2017jxz}. Analytical continuation is implemented by introducing an infinitesimal positive imaginary part to $\hat{s}$ across physical singularities in master integrals, which  correspond to $\sqrt{\hat{s}}=2m_H,m_W+m_t,2m_t,2m_t+m_Z$ or $2m_t+m_H$ when intermediate particles go on-shell. Renormalization employs the on-shell scheme for masses and fields, while the electromagnetic coupling $\alpha$ is renormalized in the $G_{\mu}$ scheme \cite{Denner:2019vbn}.


The cross section can be obtained after integrating over phase space,
\begin{align}\label{eq:Xsection}
\sigma^{\rm LO(NLO)} &= \frac{1}{512\pi}\int_0^1 \mathrm{d}x_{1}\int_0^1 \mathrm{d}x_{2} \int^{\hat{t}_{+}}_{\hat{t}_{-}}\mathrm{d}\hat{t} \nonumber \\
& \times \frac{1}{\hat{s}^{2}}f_{g/p}(x_{1},\mu)f_{g/p}(x_{2},\mu) M^{\rm LO(NLO)},
\end{align}
where $\hat{t}^{\pm}=m^{2}_{H}-\frac{\hat{s}}{2}(1\mp\sqrt{1-{4m^{2}_{H}}/{\hat{s}}})$, $f_{g/p}(x,\mu)$ denotes the gluon distribution function inside of the proton, $\mu$ represents the factorization scale, and $M^{\rm LO(NLO)}$ are given by
\begin{align}
M^{\rm LO}&=|F_{1}^{(0)}|^2+|F_{2}^{(0)}|^2 \;,\\
M^{\rm NLO}&=|F_{1}^{(0)}+F_{1}^{(1)}|^2-|F_{1}^{(1)}|^2\nonumber\\
&+|F_{2}^{(0)}+F_{2}^{(1)}|^2  -|F_{2}^{(1)}|^2,
\end{align}
where $F_i^{(0)}$ and $F_i^{(1)}$ correspond to the lowest order and the next-order terms in the $\alpha$ expansion of the form factors. The phase-space integration is carried out by optimized sampling techniques implemented through the {\tt Parni} package \cite{vanHameren:2007pt}.

Masses are set as
\begin{eqnarray}
\frac{m_H^2}{m_t^2}=\frac{12}{23},~\frac{m_Z^2}{m_t^2}=\frac{23}{83},~\frac{m_W^2}{m_t^2}=\frac{14}{65},
\end{eqnarray}
with $m_t=172.69$  GeV \cite{ParticleDataGroup:2022pth}. The electromagnetic coupling $\alpha=7.512\times 10^{-3}$ is derived from the Fermi constant $G_{\mu}=1.166378\times 10^{-5}~{\rm GeV}^{-2}$ via:
\begin{eqnarray}
\alpha=\frac{\sqrt{2}}{\pi}G_{\mu}m_W^2 \left(1-\frac{m_W^2}{m_Z^2}\right).
 \end{eqnarray}
The Cabibbo-Kobayashi-Maskawa mixing matrix is set to diagonal. NNPDF3.1\cite{NNPDF:2017mvq} serves as the default PDF set, with $\rm NNPDF31\underline{~}nlo\underline{~}as\underline{~}0118$ for both LO and NLO calculations. The running of strong coupling $\alpha_s$ with two-loop accuracy is provided by the {\tt LHAPDF6} library \cite{Buckley:2014ana} including five active flavours. Default renormalization and factorization scales are $\mu=M_{HH}/2$, where $M_{HH}$ is the invariant mass of the produced Higgs pair. 

We generate $1.8\times 10^4$ events at LO and then these events are reweighted to NLO. This enables us calculate the $\cal{K}$ factors for total and differencial cross sections. To  reduce statistical uncertainties for the the $M_{HH}$ and $p_T$ distribution, we compute an additional 400 reweighted events per bin.

\begin{table}[htbp]
\begin{center}
\renewcommand\arraystretch{1.8}
\begin{tabular}{ccccccc}
\hline
\hline
$\mu$  & $M_{HH}/2$  & $\sqrt{p_{T}^2+m_H^2}$ & $m_H$ \\
\hline
LO   & 19.96(6)   & 21.11(7)  & 25.09(8) \\
NLO  & 19.12(6)   & 20.21(6)  & 23.94(8) \\
${\cal K}$-factor  & 0.958(1)   &     0.957(1)   &   0.954(1)     \\
\hline
\hline
\end{tabular}
\caption{LO and NLO ${\rm EW}$ corrected integrated cross sections (in fb) with $\sqrt{s}= 14$ TeV based on $1.8\times 10^4$ reweighted events. The uncertainties arise from statistical errors in phase space integration.}
\label{total-cs}
\end{center}
\end{table}

Tab. \ref{total-cs} shows LO and NLO EW corrected total cross sections at $\sqrt{s}=14$ TeV for varying $\mu$. For both LO and NLO cross sections, there are around $20\%$ differences arising from $\alpha_s$ and PDF sensitivity,  mitigable via higher-order QCD corrections \cite{Jones:2023uzh}. In contrast, the $\cal{K}$ factor ranges from 0.954 to 0.958 and remains stable, indicating that the EW corrections approximately factorize from the QCD corrections. The $\cal{K}$ factor also suggest that EW corrections beyond NLO contribute at the per-mille level for the total cross section.



\begin{figure}[htbp]
\begin{center}
    \includegraphics[width=0.6\textwidth]{figures/NLO_FullEW_mhh.eps}
\caption{\label{mhh}
Invariant mass distribution of the Higgs pair  with $\sqrt{s}= 14$ TeV. The upper plot shows absolute predictions, and the lower panel displays the differential ${\cal K}$-factor with error bars representing statistical errors.}
\end{center}
\end{figure}

\begin{figure}[htbp]
\begin{center}
    \includegraphics[width=0.6\textwidth]{figures/NLO_FullEW_pth.eps}
\caption{\label{pth}
Transverse momentum distribution of one of the two Higgs bosons  with $\sqrt{s}= 14$ TeV. The upper plot shows absolute predictions, and the lower panel displays the differential ${\cal K}$-factor with error bars representing statistical errors.}
\end{center}
\end{figure}

In Fig. \ref{mhh}, we present the $M_{HH}$ distribution. A significant positive correction of approximately $+15\%$ is observed in the first bin. In fact, we find that the EW correction for phase space points near the $HH$ production threshold can exceed $+70\%$. As $M_{HH}$ increases, the $\mathcal{K}$ factor initially decreases dramatically and then moderates farther from the threshold. A similar pattern occurs for the $p_{T}$ distribution in Fig. \ref{pth}, where the correction is initially positive before turning negative around 100 GeV. For regions of large $M_{HH}$ or large $p_{T}$, the NLO EW correction is approximately $-10\%$. Note that, while the corrections can reach -$30\%$ for the squared matrix element at $\sqrt{\hat{s}} \approx 14\ {\rm TeV}$, the extreme suppression of gluon luminosity at high energy makes these impact insignificant. In Fig. \ref{yh}, we display the rapidity distribution of one Higgs boson. The $\mathcal{K}$ factor is nearly flat, approximately $0.96$, matching that of the total cross section.

\begin{figure}[htbp]
\begin{center}
    \includegraphics[width=0.6\textwidth]{figures/NLO_FullEW_yh.eps}
\caption{\label{yh} 
Rapidity distribution of one of the two Higgs bosons  with $\sqrt{s}= 14$ TeV. The upper plot shows absolute predictions, and the lower panel displays the differential ${\cal K}$-factor with error bars representing statistical errors.}
\end{center}
\end{figure}